\section{Introduction to SNARKs}

Peggy has done some very difficult calculation. She wants to prove to
Victor that she did it. Victor wants to check that Peggy did it, but he
is too lazy to redo the whole calculation himself.

\begin{itemize}
\item
  Maybe Peggy wants to keep part of the calculation secret. Maybe her
  calculation was ``find a solution to this puzzle,'' and she wants to
  prove that she found a solution without saying what the solution is.
\item
  Maybe it is just a really long, annoying calculation, and Victor does
  not have the energy to check it all line-by-line.
\end{itemize}

A \emph{SNARK}  lets Peggy (the ``prover'') send Victor (the
``verifier'') a short proof that she has indeed done the calculation
correctly. The proof will be much shorter than the original calculation,
and Victor's verification is much faster. (As a tradeoff, writing a
SNARK proof of a calculation is much slower than just doing the
calculation.)

The name stands for:

\begin{itemize}
\item
  \emph{Succinct}: the proof length is short (in fact, it is a constant
  length, independent of how long the problem is).
\item
  \emph{Non-interactive}: the protocol does not require back-and-forth
  communication.
\item
  \emph{Argument}: basically a proof. There is a technical difference,
  but we will not worry about it.
\item
  \emph{of Knowledge}: the proof does not just show the system of
  equations has a solution; it also shows that the prover knows one.
\end{itemize}

We will not discuss it here, but it is also possible and frequently
useful to make a \emph{zero knowledge (zk)} SNARK. These are typically
called ``zkSNARKs.'' This gives Peggy a guarantee that Victor will not
learn anything about the intermediate steps in her calculation, aside
from any particular steps Peggy chooses to reveal.

\subsection{What can you do with a SNARK?}

One answer: You can prove that you have a solution to a system of
equations. Sounds pretty boring, unless you are an algebra student.

Slightly better answer: You can prove that you have executed a program
correctly, revealing some or all of the inputs and outputs, as you
please. For example: You know a message \(M\) such that
\(\operatorname{sha}(M) = \ 0xa91af3ac...\), but you do not want to
reveal \(M\). Or: You only want to reveal the first 30 bytes of \(M\).
Or: You know a message \(M\), and a digital signature proving that \(M\)
was signed by {[}trusted authority{]}, such that a certain neural
network, run on the input \(M\), outputs ``Good.''

One recent application along these lines is TLSNotary.\footnote{https://tlsnotary.org}
TLSNotary lets you certify a transcript of communications with a server
in a privacy-preserving way: You only reveal the parts you want to.
